\chapter*{Conclusion}
\addcontentsline{toc}{chapter}{Conclusion}

L'échantillon étudié appartient à une classe de menaces encore peu documentée : les
codes malveillants qui interrogent un modèle de langage pendant leur propre
exécution. Le caractériser a demandé de reconstruire un pipeline dont presque rien
ne se répétait d'une exécution à l'autre. L'analyse statique a livré six tâches
codées en dur, une clé de session en clair, un chemin de compilation oublié et les
prompts système embarqués dans le binaire ; l'analyse dynamique a confirmé les noms
de tâches, les artefacts produits, l'endpoint d'exfiltration et un délai
d'expiration côté client d'environ 170 secondes. J'en ai tiré un jeu de règles YARA,
cinq règles Sigma et les fondations d'un plugin Volatility.

Chez PromptLock, la génération varie mais l'orchestrateur non. Le code Lua diffère à
chaque exécution, la note de rançon aussi, le contenu chiffré ne se répète jamais.
L'enveloppe qui pilote l'ensemble, elle, ne bouge pas : imports Go, prompts système,
noms de tâches, chemin de compilation, format des échanges avec l'API. La détection
reste donc possible, à condition de viser la structure qui produit le contenu plutôt
que le contenu lui-même.

Cette stabilité n'a rien d'acquis. PromptFlux, documenté par le GTIG en novembre
2025, vise l'inverse : faire réécrire son propre code source par le modèle. La
tentative n'aboutit pas, mais l'échantillon analysé ici en est sans doute l'état
antérieur.

Ni ATT\&CK ni ATLAS ne décrivent l'usage d'un modèle comme outil au sein d'un code
malveillant, par opposition à l'attaque menée contre un système d'IA. Le GTIG s'est
lui-même rabattu sur des sous-techniques existantes pour cartographier PromptFlux.
Tant que ce vocabulaire manque, deux analystes qui observent le même comportement le
nomment autrement, et le renseignement circule mal.

Un seul échantillon, donc, et probablement un prototype de recherche : adresse
Bitcoin de test, nom de compilation conservé. Les règles n'ont pas été confrontées à
des variantes modifiées pour les contourner. Leur taux de faux positifs reste à
mesurer, et sans cette mesure l'équipe ne reprendra rien : une règle qui bruite sur
un parc entier finit désactivée, juste ou pas.

Un binaire qui chargerait ses instructions depuis un serveur, ou les dériverait d'un
modèle local, ne laisserait plus aucune prise à YARA et reporterait l'effort sur la
télémétrie comportementale. S'y ajoute la baisse du coût d'entrée : le marché des
outils offensifs assistés par IA s'est structuré en 2025, et ce qui demande
aujourd'hui de savoir intégrer une API dans un binaire Go sera bientôt vendu clé en
main.

Les règles écrites ici vieilliront, plus vite que celles d'un rançongiciel
classique. La question qui les a produites se repose à chaque nouvel échantillon :
qu'est-ce qui ne peut pas varier sans que la menace cesse de fonctionner ?